% ---
% title: C++17 结构化绑定
% date: 2026-07-12
% author: Crystal-Sky
% tags: C++, C++17, 基础语法, 结构化绑定
% summary: 系统介绍 C++17 结构化绑定的基本语法、数组与结构体拆分、pair/tuple、引用绑定及常见使用场景。
% slug: cpp-structured-bindings
% course: C++
% course_slug: cpp
% chapter: 基础语法
% chapter_order: 1
% lesson_order: 1
% lesson_type: 基础语法
% ---

\documentclass[UTF8,12pt,a4paper]{ctexart}

\usepackage{geometry}
\usepackage{xcolor}
\usepackage{listings}
\usepackage{amsmath}
\usepackage{amssymb}
\usepackage{booktabs}
\usepackage{enumitem}
\usepackage{hyperref}
\usepackage{tcolorbox}
\usepackage{fancyhdr}

\geometry{left=2.4cm,right=2.4cm,top=2.5cm,bottom=2.5cm}

\hypersetup{
    colorlinks=true,
    linkcolor=blue,
    urlcolor=blue
}

\pagestyle{fancy}
\fancyhf{}
\fancyhead[L]{C++17 结构化绑定}
\fancyfoot[C]{\thepage}

\definecolor{codebg}{RGB}{247,248,250}
\definecolor{keywordcolor}{RGB}{0,92,197}
\definecolor{commentcolor}{RGB}{0,128,0}
\definecolor{stringcolor}{RGB}{163,21,21}

\lstset{
    language=C++,
    basicstyle=\ttfamily\small,
    keywordstyle=\color{keywordcolor}\bfseries,
    commentstyle=\color{commentcolor},
    stringstyle=\color{stringcolor},
    backgroundcolor=\color{codebg},
    frame=single,
    numbers=left,
    numberstyle=\tiny\color{gray},
    stepnumber=1,
    numbersep=8pt,
    tabsize=4,
    showstringspaces=false,
    breaklines=true,
    columns=fullflexible,
    keepspaces=true
}

\newtcolorbox{conceptbox}[1]{
    colback=blue!3,
    colframe=blue!60!black,
    title=#1,
    fonttitle=\bfseries
}

\newtcolorbox{warningbox}[1]{
    colback=red!3,
    colframe=red!65!black,
    title=#1,
    fonttitle=\bfseries
}

\newtcolorbox{examplebox}[1]{
    colback=green!3,
    colframe=green!50!black,
    title=#1,
    fonttitle=\bfseries
}

\title{\textbf{C++17 结构化绑定}}
\author{}
\date{}

\begin{document}

\maketitle

\section{学习目标}

完成本节学习后，应能够：

\begin{enumerate}[label=\arabic*.]
    \item 理解结构化绑定的作用和基本语法；
    \item 使用结构化绑定拆分数组、结构体、\texttt{pair} 和 \texttt{tuple}；
    \item 正确区分值绑定、引用绑定和常量引用绑定；
    \item 在范围 \texttt{for} 循环和 STL 容器中使用结构化绑定；
    \item 识别结构化绑定中的常见错误。
\end{enumerate}

\section{什么是结构化绑定}

结构化绑定（Structured Bindings）是 C++17 引入的一种语法。

它允许程序员将一个由多个元素组成的对象拆分为若干个变量，从而减少使用
\texttt{first}、\texttt{second}、\texttt{get<0>()} 等写法。

\begin{conceptbox}{基本思想}
结构化绑定可以理解为：

\[
\text{一个复合对象}
\quad \longrightarrow \quad
\text{多个具有名字的变量}
\]

例如，一个坐标对象中包含横坐标和纵坐标，可以直接拆分为变量 \texttt{x} 和 \texttt{y}。
\end{conceptbox}

传统写法：

\begin{lstlisting}
std::pair<int, int> point = {10, 20};

int x = point.first;
int y = point.second;
\end{lstlisting}

使用结构化绑定：

\begin{lstlisting}
std::pair<int, int> point = {10, 20};

auto [x, y] = point;
\end{lstlisting}

此时：

\begin{itemize}
    \item \texttt{x} 的值为 10；
    \item \texttt{y} 的值为 20。
\end{itemize}

\section{基本语法}

结构化绑定的一般形式如下：

\begin{lstlisting}
auto [变量1, 变量2, ...] = 表达式;
\end{lstlisting}

也可以使用引用或常量引用：

\begin{lstlisting}
auto [a, b] = object;          // 值绑定
auto& [a, b] = object;         // 引用绑定
const auto& [a, b] = object;   // 常量引用绑定
\end{lstlisting}

\subsection{变量数量必须匹配}

结构化绑定左侧变量的数量必须与对象中可拆分元素的数量一致。

正确示例：

\begin{lstlisting}
std::pair<int, int> p = {1, 2};
auto [x, y] = p;
\end{lstlisting}

错误示例：

\begin{lstlisting}
std::pair<int, int> p = {1, 2};
auto [x] = p;       // 错误：pair 中有两个元素
\end{lstlisting}

\section{绑定数组}

结构化绑定可以拆分普通数组。

\begin{lstlisting}
#include <iostream>
using namespace std;

int main() {
    int arr[3] = {10, 20, 30};

    auto [a, b, c] = arr;

    cout << a << ' ' << b << ' ' << c << '\n';
    return 0;
}
\end{lstlisting}

输出结果：

\begin{verbatim}
10 20 30
\end{verbatim}

这里的 \texttt{a}、\texttt{b}、\texttt{c} 是数组元素的副本。

修改这些变量不会改变原数组。

\begin{lstlisting}
int arr[3] = {10, 20, 30};

auto [a, b, c] = arr;
a = 100;

cout << arr[0] << '\n';  // 输出 10
\end{lstlisting}

如果希望通过绑定变量修改原数组，应使用引用绑定：

\begin{lstlisting}
int arr[3] = {10, 20, 30};

auto& [a, b, c] = arr;
a = 100;

cout << arr[0] << '\n';  // 输出 100
\end{lstlisting}

\section{绑定结构体}

结构化绑定可以拆分成员均为公开成员的简单结构体。

\begin{lstlisting}
#include <iostream>
#include <string>
using namespace std;

struct Student {
    string name;
    int age;
    double score;
};

int main() {
    Student stu{"Alice", 20, 95.5};

    auto [name, age, score] = stu;

    cout << name << '\n';
    cout << age << '\n';
    cout << score << '\n';

    return 0;
}
\end{lstlisting}

结构体成员的绑定顺序与成员的声明顺序一致：

\begin{lstlisting}
struct Student {
    string name;   // 第一个
    int age;       // 第二个
    double score;  // 第三个
};
\end{lstlisting}

因此：

\begin{lstlisting}
auto [name, age, score] = stu;
\end{lstlisting}

分别对应 \texttt{stu.name}、\texttt{stu.age} 和 \texttt{stu.score}。

\begin{warningbox}{注意}
结构化绑定依赖成员声明顺序，而不是变量名称。

左侧变量可以使用任意合法名称，但顺序不能写错。
\end{warningbox}

例如：

\begin{lstlisting}
auto [a, b, c] = stu;
\end{lstlisting}

其中：

\begin{itemize}
    \item \texttt{a} 对应 \texttt{name}；
    \item \texttt{b} 对应 \texttt{age}；
    \item \texttt{c} 对应 \texttt{score}。
\end{itemize}

\section{绑定 \texttt{std::pair}}

\texttt{std::pair} 中包含两个元素，分别是 \texttt{first} 和 \texttt{second}。

传统写法：

\begin{lstlisting}
std::pair<std::string, int> student = {"Tom", 18};

std::string name = student.first;
int age = student.second;
\end{lstlisting}

结构化绑定写法：

\begin{lstlisting}
std::pair<std::string, int> student = {"Tom", 18};

auto [name, age] = student;
\end{lstlisting}

完整示例：

\begin{lstlisting}
#include <iostream>
#include <string>
#include <utility>
using namespace std;

int main() {
    pair<string, int> student = {"Tom", 18};

    auto [name, age] = student;

    cout << "姓名：" << name << '\n';
    cout << "年龄：" << age << '\n';

    return 0;
}
\end{lstlisting}

\section{绑定 \texttt{std::tuple}}

\texttt{std::tuple} 可以存储任意数量、任意类型的元素。

传统写法：

\begin{lstlisting}
std::tuple<std::string, int, double> stu{"Bob", 19, 88.5};

std::string name = std::get<0>(stu);
int age = std::get<1>(stu);
double score = std::get<2>(stu);
\end{lstlisting}

结构化绑定写法：

\begin{lstlisting}
auto [name, age, score] = stu;
\end{lstlisting}

完整示例：

\begin{lstlisting}
#include <iostream>
#include <string>
#include <tuple>
using namespace std;

int main() {
    tuple<string, int, double> stu{"Bob", 19, 88.5};

    auto [name, age, score] = stu;

    cout << name << ' ' << age << ' ' << score << '\n';

    return 0;
}
\end{lstlisting}

\section{值绑定与引用绑定}

结构化绑定中最重要的内容之一，是区分值绑定和引用绑定。

\subsection{值绑定}

语法：

\begin{lstlisting}
auto [a, b] = object;
\end{lstlisting}

此时会根据原对象创建一份新的绑定对象。

通过绑定变量进行修改，通常不会影响原对象。

\begin{lstlisting}
std::pair<int, int> p = {10, 20};

auto [x, y] = p;
x = 100;

cout << p.first << '\n';  // 输出 10
\end{lstlisting}

\subsection{引用绑定}

语法：

\begin{lstlisting}
auto& [a, b] = object;
\end{lstlisting}

绑定变量直接引用原对象中的元素。

\begin{lstlisting}
std::pair<int, int> p = {10, 20};

auto& [x, y] = p;
x = 100;

cout << p.first << '\n';  // 输出 100
\end{lstlisting}

\subsection{常量引用绑定}

语法：

\begin{lstlisting}
const auto& [a, b] = object;
\end{lstlisting}

这种方式不会复制对象，同时禁止通过绑定变量修改元素。

\begin{lstlisting}
std::pair<int, int> p = {10, 20};

const auto& [x, y] = p;

cout << x << ' ' << y << '\n';
// x = 100;  // 错误：x 是只读的
\end{lstlisting}

\subsection{三种方式对比}

\begin{center}
\begin{tabular}{lll}
\toprule
写法 & 是否复制 & 能否修改原对象 \\
\midrule
\texttt{auto [a,b]} & 是 & 否 \\
\texttt{auto\& [a,b]} & 否 & 是 \\
\texttt{const auto\& [a,b]} & 否 & 否 \\
\bottomrule
\end{tabular}
\end{center}

\begin{conceptbox}{选择建议}
\begin{itemize}
    \item 对象较小，且不需要修改原对象时，可使用 \texttt{auto}；
    \item 需要修改原对象时，使用 \texttt{auto\&}；
    \item 对象较大，只读访问时，优先使用 \texttt{const auto\&}。
\end{itemize}
\end{conceptbox}

\section{在函数返回值中使用}

结构化绑定非常适合接收包含多个结果的函数返回值。

\subsection{返回 \texttt{pair}}

\begin{lstlisting}
#include <iostream>
#include <utility>
using namespace std;

pair<int, int> divide(int a, int b) {
    return {a / b, a % b};
}

int main() {
    auto [quotient, remainder] = divide(17, 5);

    cout << "商：" << quotient << '\n';
    cout << "余数：" << remainder << '\n';

    return 0;
}
\end{lstlisting}

输出：

\begin{verbatim}
商：3
余数：2
\end{verbatim}

\subsection{返回 \texttt{tuple}}

\begin{lstlisting}
#include <iostream>
#include <tuple>
using namespace std;

tuple<int, int, double> analyze(int a, int b) {
    int sum = a + b;
    int product = a * b;
    double average = (a + b) / 2.0;

    return {sum, product, average};
}

int main() {
    auto [sum, product, average] = analyze(4, 6);

    cout << sum << '\n';
    cout << product << '\n';
    cout << average << '\n';

    return 0;
}
\end{lstlisting}

这种写法比使用多个输出参数更加直观。

\section{在范围 \texttt{for} 循环中使用}

结构化绑定经常与范围 \texttt{for} 循环结合使用。

\subsection{遍历 \texttt{map}}

传统写法：

\begin{lstlisting}
for (const auto& item : mp) {
    cout << item.first << ' ' << item.second << '\n';
}
\end{lstlisting}

结构化绑定写法：

\begin{lstlisting}
for (const auto& [key, value] : mp) {
    cout << key << ' ' << value << '\n';
}
\end{lstlisting}

完整示例：

\begin{lstlisting}
#include <iostream>
#include <map>
#include <string>
using namespace std;

int main() {
    map<string, int> scores{
        {"Alice", 95},
        {"Bob", 88},
        {"Tom", 91}
    };

    for (const auto& [name, score] : scores) {
        cout << name << ": " << score << '\n';
    }

    return 0;
}
\end{lstlisting}

\subsection{修改 \texttt{map} 中的值}

如果需要修改容器中的值，应使用引用：

\begin{lstlisting}
for (auto& [name, score] : scores) {
    score += 5;
}
\end{lstlisting}

此时每个学生的分数都会增加 5。

\begin{warningbox}{注意}
\texttt{std::map} 的键是只读的。

即使使用 \texttt{auto\& [key, value]}，也不能修改 \texttt{key}，但可以修改 \texttt{value}。
\end{warningbox}

\section{与集合插入操作结合}

许多 STL 操作会返回 \texttt{pair}，结构化绑定可以使代码更清晰。

例如，\texttt{set::insert} 会返回：

\begin{lstlisting}
pair<iterator, bool>
\end{lstlisting}

其中：

\begin{itemize}
    \item 第一个元素是指向对应元素的迭代器；
    \item 第二个元素表示是否成功插入。
\end{itemize}

示例：

\begin{lstlisting}
#include <iostream>
#include <set>
using namespace std;

int main() {
    set<int> numbers;

    auto [it1, success1] = numbers.insert(10);
    auto [it2, success2] = numbers.insert(10);

    cout << boolalpha;
    cout << success1 << '\n';  // true
    cout << success2 << '\n';  // false

    return 0;
}
\end{lstlisting}

这种写法比下面的传统形式更容易理解：

\begin{lstlisting}
auto result = numbers.insert(10);

auto it = result.first;
bool success = result.second;
\end{lstlisting}

\section{结构化绑定的作用域}

结构化绑定变量与普通局部变量一样，受到作用域限制。

\begin{lstlisting}
int main() {
    {
        std::pair<int, int> p = {1, 2};
        auto [x, y] = p;

        std::cout << x << ' ' << y << '\n';
    }

    // std::cout << x;  // 错误：x 已离开作用域
}
\end{lstlisting}

结构化绑定变量也不能与当前作用域内已有变量重名。

\begin{lstlisting}
int x = 10;
std::pair<int, int> p = {1, 2};

auto [x, y] = p;  // 错误：x 重复定义
\end{lstlisting}

\section{常见错误}

\subsection{忘记启用 C++17}

结构化绑定属于 C++17 特性。

使用 GCC 或 Clang 编译时，应指定：

\begin{verbatim}
g++ main.cpp -std=c++17 -o main
\end{verbatim}

如果未启用 C++17，可能出现类似错误：

\begin{verbatim}
structured bindings only available with -std=c++17
\end{verbatim}

\subsection{绑定变量数量不匹配}

\begin{lstlisting}
std::tuple<int, int, int> t{1, 2, 3};

auto [a, b] = t;  // 错误：tuple 中有三个元素
\end{lstlisting}

应改为：

\begin{lstlisting}
auto [a, b, c] = t;
\end{lstlisting}

\subsection{误以为值绑定会修改原对象}

\begin{lstlisting}
std::pair<int, int> p{1, 2};

auto [x, y] = p;
x = 100;

cout << p.first;  // 仍然是 1
\end{lstlisting}

需要修改原对象时应使用：

\begin{lstlisting}
auto& [x, y] = p;
\end{lstlisting}

\subsection{在循环中不必要地复制}

下面的代码会在每次循环时复制元素：

\begin{lstlisting}
for (auto [key, value] : mp) {
    cout << key << ' ' << value << '\n';
}
\end{lstlisting}

只读遍历时更推荐：

\begin{lstlisting}
for (const auto& [key, value] : mp) {
    cout << key << ' ' << value << '\n';
}
\end{lstlisting}

\subsection{尝试绑定私有成员}

结构化绑定不能直接拆分具有不可访问成员的普通类对象。

\begin{lstlisting}
class Student {
private:
    string name;
    int age;
};
\end{lstlisting}

直接写：

\begin{lstlisting}
Student stu;
auto [name, age] = stu;
\end{lstlisting}

通常会因为成员不可访问而编译失败。

在基础教学中，可以先将结构化绑定用于简单结构体、数组、\texttt{pair} 和 \texttt{tuple}。

\section{综合示例：学生成绩统计}

下面的程序使用结构体、\texttt{vector}、结构化绑定和函数返回值，完成学生成绩统计。

\begin{lstlisting}
#include <iostream>
#include <string>
#include <tuple>
#include <vector>
using namespace std;

struct Student {
    string name;
    int score;
};

tuple<int, int, double> analyze(const vector<Student>& students) {
    int maximum = students.front().score;
    int minimum = students.front().score;
    int sum = 0;

    for (const auto& [name, score] : students) {
        if (score > maximum) {
            maximum = score;
        }

        if (score < minimum) {
            minimum = score;
        }

        sum += score;
    }

    double average =
        static_cast<double>(sum) / students.size();

    return {maximum, minimum, average};
}

int main() {
    vector<Student> students{
        {"Alice", 95},
        {"Bob", 82},
        {"Tom", 91}
    };

    for (const auto& [name, score] : students) {
        cout << name << ": " << score << '\n';
    }

    auto [maximum, minimum, average] = analyze(students);

    cout << "最高分：" << maximum << '\n';
    cout << "最低分：" << minimum << '\n';
    cout << "平均分：" << average << '\n';

    return 0;
}
\end{lstlisting}


\section{知识总结}

\begin{conceptbox}{本节核心内容}
\begin{enumerate}[label=\arabic*.]
    \item 结构化绑定是 C++17 引入的语法；
    \item 基本形式为 \texttt{auto [a,b,...] = object;}；
    \item 可以绑定数组、简单结构体、\texttt{pair} 和 \texttt{tuple}；
    \item \texttt{auto} 表示值绑定，\texttt{auto\&} 表示引用绑定；
    \item \texttt{const auto\&} 适合高效的只读访问；
    \item 结构化绑定常用于函数多返回值和 STL 容器遍历；
    \item 编译时需要开启 C++17 标准。
\end{enumerate}
\end{conceptbox}

\section{板书式速记}

\begin{lstlisting}
// 1. 基本语法
auto [a, b] = object;

// 2. 修改原对象
auto& [a, b] = object;

// 3. 只读且避免复制
const auto& [a, b] = object;

// 4. 遍历 map
for (const auto& [key, value] : mp) {
    // ...
}

// 5. 接收函数的多个返回值
auto [x, y, z] = function();
\end{lstlisting}

\end{document}
